\section{Finite modules and the exhaustion interface}\label{sec:setup}

Our algebraic calculations need only countability and ICC. The geometric
class enters at one precise point: the exhaustion statement below says
that the explicit models contain every bifinite correspondence. We first
fix the norm convention for their construction.

For completeness, $L_\mu(\Gamma)$ is a factor when $\Gamma$ is ICC:
the absolute values of the Fourier coefficients of a central $L^2$ vector
are constant on conjugacy classes, so square summability kills every
nonidentity coefficient.

\subsection{Relative tensor products and elementary dimensions}
All Hilbert-space inner products are linear in the second variable.
A vector $\xi$ in an $M$--$N$ correspondence is \emph{right bounded}
if $R_\xi\widehat a=\xi a$ extends to a bounded map from $L^2(N)$.
For an $N$--$P$ correspondence, the relative tensor product is the
completion, after quotienting out zero-norm vectors, of the tensors
with right-bounded first vector and inner product
\begin{equation}\label{eq:relative-inner-product}
 \langle\xi\otimes_N\eta,\xi'\otimes_N\eta'\rangle
 =\langle\eta,(R_\xi^*R_{\xi'})\eta'\rangle.
\end{equation}
The right-module map $R_\xi^*R_{\xi'}$ is left multiplication by an
element of $N$. Balancing means
$(\xi a)\otimes_N\eta=\xi\otimes_N(a\eta)$ for $a\in N$.
This is the finite tracial form of Connes's relative tensor product
\cite[Chapter V, Appendix B, Theorems~13 and~20]{Connes1994}.

A finite right $N$-module has a column realization $pL^2(N)^n$,
with dimension $(\Tr_n\otimes\tau_N)(p)$. The vectors in $pN^n$
are right bounded and dense. The projection model and its dimension
are independent of the chosen realization
\cite[Theorem~10.1.1 and Proposition~10.1.3]{Jones2009}.
We prove the normality of the graph models directly in
Lemmas~\ref{lem:calculus-induction} and~\ref{lem:calculus-graph-model}.
The same elementary construction precedes the geometric argument in the
companion, so its exhaustion proof has no dependence on our all-maps or
summand classification. For $n=\dim\sigma$, the elementary dimensions are
$n[\Lambda:B]$ and $n[\Gamma:A]$ in left/right order. The dimensions of
\emph{actual summands} will require the full-stabilizer calculation.

\subsection{The geometric input}
The real simple factors in Definition~\ref{def:class} have higher rank
or are locally $\Sp(n,1)$, $n\geq2$, or $F_{4(-20)}$. At a finite
place their relative rank is at least two. Complex places are treated as
real Lie groups in differential arguments.
Property~$(T)$ passes between a locally compact group and a lattice and
between a discrete group and a finite-index subgroup; see \cite{BHV2008}. In particular,
that class is closed under abstract commensurability among countably
infinite ICC groups. The geometric companion treats arbitrary lattices
in the permitted finite products, without a dense-projection hypothesis.
We use exactly the following theorem.

\begin{proposition}[Arithmetic exhaustion, companion Theorem~1.2]
\label{prop:recovery-exhaustion}
Let $\Gamma\in\mathscr K$, let $\Lambda$ be any countably infinite ICC
group, and let $\mu,\omega$ be arbitrary normalized scalar cocycles.
Every nonzero separable bifinite
$L_\mu(\Gamma)$--$L_\omega(\Lambda)$ correspondence is a closed bimodule
summand of a finite direct sum of the models
$E(A,B,\delta,\sigma)$ in Equation~\eqref{eq:intro-model}, with the
exact multiplier in Equation~\eqref{eq:intro-multiplier}.
\end{proposition}

This is the Arithmetic exhaustion theorem of the complete companion
\cite[Theorem~1.2]{ArithmeticCompanion}. Its conclusion
retains actual finite-index subgroup isomorphisms and their exact scalar
ratios. We do not assume that the projective representations have finite
image, that either cocycle is finite type, or that the original
correspondence is irreducible. The algebraic proofs that follow are
independent of this proposition until its explicit application at the end
of Section~\ref{sec:calculus}.
